\cotefile{136}
\scanpage{189}\mspage{- 64 -}\hfill IX

une propriété nettement plus forte que celle énoncée dans 5.7, qui pourra alors servir de caractérisation intrinsèque de la partie fixe de $T_p(A_K)$ (en termes de la seule structure de pro-groupe de Barsotti-Tate de $T_p(A_K)$), - du moins dans le cas d'inégales caractéristiques :

\uline{Théorème} 5.8. (M. RAYNAUD). \uline{Soit} $A_K$ \uline{un schéma abélien à réduction semi-stable sur le corps des fonctions} $K$ \uline{d'un trait hensélien} $S$. \uline{Soit} $\handwritten{\ell}$ \uline{un nombre premier} (\uline{le cas intéressant étant celui où} $\handwritten{\ell}=p$), \uline{et considérons le pro-}$\handwritten{\ell}$\uline{-groupe de Barsotti-Tate} $X_{\handwritten{\eta}}=T_{\handwritten{\ell}}(A_K)$ \uline{et sa} ``\uline{partie fixe}'' $X_{\handwritten{\eta}}^f$, \uline{qui est la fibre générique du pro-}$\handwritten{\ell}$\uline{-groupe de Barsotti-Tate} $X^f=T_{\handwritten{\ell}}(A^{\unsure{o|0}})^f$. \uline{Supposons}, \uline{si} $\handwritten{\ell}=p$, \uline{qu'on ait} $\mathrm{car.}\ K=0$ (\uline{de sorte que nous disposons sur} $S$ \uline{du théorème de TATE} [33, th.4]). \uline{Soit} $F$ \uline{un pro-}$\handwritten{\ell}$\uline{-groupe de Barsotti-Tate sur} $S$. \uline{Alors tout homomorphisme} $F_{\handwritten{\eta}}\handwritten{\longrightarrow}X_{\handwritten{\eta}}$ \uline{se factorise par} $X_{\handwritten{\eta}}^f$.

Passant aux ind-$\handwritten{\ell}$-groupes associés aux pro-$\handwritten{\ell}$-groupes de Barsotti-Tate envisagés, l'énoncé précédent devient équivalent au suivant, grâce au théorème cité de TATE :

\uline{Corollaire} 5.9. \uline{Si} $F$ \uline{est un} ind-$\handwritten{\ell}$\uline{-groupe de Barsotti-Tate}, \uline{alors tout homomorphisme} $u_{\handwritten{\eta}}: F_{\handwritten{\eta}}\handwritten{\longrightarrow} A_{\handwritten{\eta}}$ \uline{se prolonge} (\uline{de façon unique}) \uline{en un homomorphisme} $u: F\handwritten{\longrightarrow} A$.

En d'autres termes, la propriété ``néronienne'' caractéristique (1.1.2) du modèle de Néron reste valable pour des homomorphismes de ind-schémas en groupes de Barsotti-Tate, sans hypothèse de lissité sur celui-ci. (Seul le cas des groupes ind-étales résulterait directement de la définition.)

% archive sheet number: 188
% archive numbering: 382 (pencil and printed label, foot).
